% Table 3 -- finetuned MSR-VTT. Both methods given both finetuning learning rates.
\begin{table}[t]
\centering
\caption{Fine-tuned retrieval on MSR-VTT, R@1 (ITM protocol). Every row starts from a
150k-budget pretrain and uses identical finetuning data and schedule. Both methods were
given both finetuning learning rates, so the comparison is best-configuration against
best-configuration: $2\!\times\!10^{-5}$ is better for both, and GRAM's volume objective
finetunes more effectively on this in-domain benchmark ($-$1.8 R@1). Adapter-only
finetuning is the binding constraint for SCA ($+$2.8 from LoRA to full finetuning); the
initialisation is not (v2a, a better zero-shot start, gains nothing). The last row adds
depth as a fifth modality with no change to the centroid computation, but it is not a
controlled comparison: it continues from the row-1 SCA finetune for a further four epochs
at $1\!\times\!10^{-4}$, so the $+$1.6 over 57.1 confounds arity with the extra pass. The
matched control -- the same checkpoint scored on T-VAS instead of T-VASD -- is pending.}
\label{tab:finetune}
\small
\setlength{\tabcolsep}{5pt}
\begin{tabular}{llccrr}
\toprule
Method & FT regime & FT lr & Init & T2V & V2T \\
\midrule
GRAM (reproduced) & full-FT & 2e-5 & GRAM & \textbf{61.2} & \textbf{60.4} \\
GRAM (reproduced) & full-FT & 1e-4 & GRAM & 58.8 & 58.2 \\
\midrule
SCA & LoRA & 2e-5 & SCA & 57.1 & 58.9 \\
SCA & LoRA & 2e-5 & SCA$_{\text{lr}}$ & 56.6 & 57.9 \\
\textbf{SCA} & full-FT & 2e-5 & SCA$_{\text{lr}}$ & \textbf{59.4} & \textbf{60.9} \\
SCA & full-FT & 1e-4 & SCA$_{\text{lr}}$ & 56.2 & 56.9 \\
\midrule
SCA (T-VASD, $k{=}5$) & LoRA & 1e-4 & SCA ft & 58.7 & 56.2 \\
\bottomrule
\end{tabular}
\end{table}
